\documentclass[10pt,a4paper]{amsart}
\usepackage[margin=19mm]{geometry}
\usepackage{amsmath,amssymb,mathtools}
\usepackage[scr=rsfs,cal=euler]{mathalfa}
\usepackage{xfrac}
\usepackage[unicode,hidelinks,bookmarks=false]{hyperref}
\newcommand{\ie}{i.e.\ }
\newcommand{\cf}{cf.\ }
\newcommand{\resp}{respectively\ }
\newcommand{\Subsection}{\subsection}
\providecommand{\nobreakdash}{\nobreak\hbox{-}}
\setlength{\emergencystretch}{2em}
\multlinegap0pt
\usepackage{bm}
\usepackage{bbm}
\usepackage{dsfont}
\usepackage{xspace}
\usepackage{cancel}
\usepackage{mathtools}
\usepackage{amsthm}
\makeatletter
\@ifundefined{theorem}{\newtheorem{theorem}{Theorem}}{}
\makeatother
\title[Onsager-Machlup Functional for SDE with Time-Varying Fractio]{Onsager-Machlup Functional for SDE with Time-Varying Fractional Noise}
\author{Yanbin Zhu, Xiaomeng Jiang,, Yong Li}
\date{Source: arXiv:2511.09300v1 (CC BY 4.0)}
\begin{document}
\maketitle

\noindent\emph{Source and licence.} Onsager-Machlup Functional for SDE with Time-Varying Fractional Noise. Version arXiv:2511.09300v1 (CC BY 4.0), distributed under the Creative Commons Attribution 4.0 International licence, as recorded in the arXiv licence metadata for this exact version and shown on the abstract page https://arxiv.org/abs/2511.09300v1. Full rights notice: licenses/arxiv-2511.09300-CC-BY-4.0.txt.\par\smallskip

\noindent\textit{Editorial scope.} Counted unit: the Theorem whose source label is \texttt{result2} in the article (source0.tex lines 1500--1512, source label \texttt{result2}), delivered together with the article's own complete original proof of it (source0.tex lines 1513--1833, one \texttt{proof} environment of 321 lines). No mathematics is written, repaired or re-derived: statement and proof are the article's own text line for line. Required context retained uncounted: fir (lines 85--87); estimate1 (lines 574--580). Every definition, display and label of those lines is retained unchanged. Omitted from this package and not redistributed: 1--84, 88--573, 581--1499, 1834--2279. Nothing inside a retained excerpt is omitted or altered: every retained body is delivered exactly as the article's lines are written, so no omission receipt of any kind accompanies this package. Citations: the retained excerpts carry no citation command at all, so no bibliography entry of the article is transcribed and no reference is redistributed. Reference adaptation: none was needed and none was made; every reference inside the delivered unit resolves inside it, so no unresolved reference or citation is produced. Printed slips kept literal: no printed word, symbol, hypothesis or number of the counted statement or of its proof was altered. Layout and class: the article's own class is \texttt{elsarticle} with packages including longtable, booktabs, amsmath, amssymb, amsfonts, amsthm, mathrsfs, graphicx, subcaption, float, afterpage, hyperref, bookmark, xcolor; the wrapper is the lane's pinned \texttt{amsart} editorial wrapper at 10pt on A4, which restates the theorem-like environments the retained text needs and adds the article's own macro definitions that the retained text needs (none), with extra wrapper packages (bm, bbm, dsfont, xspace, cancel, mathtools). The delivered numbering is the unit's own. Assets: no retained span contains a figure environment, a graphics-inclusion command, a PDF-page inclusion, a table or a TikZ picture; no image file is copied into this package, the package declares no asset dependency and no image file is redistributed with it. Licence: version arXiv:2511.09300v1 of this article is distributed under the Creative Commons Attribution 4.0 International licence, as recorded in the arXiv licence metadata for this exact version and shown on the abstract page https://arxiv.org/abs/2511.09300v1. A full rights notice is delivered at licenses/arxiv-2511.09300-CC-BY-4.0.txt. Characters: every retained excerpt is pure ASCII, so the delivered engine, XeLaTeX with its default Latin Modern text font, reports no missing character.
\begin{equation}
    X_t = x_0 + \int_0^t b_s(X_s)\,ds + \int_0^t \sigma_s\,dB^H_s, \label{fir}
\end{equation}

	\begin{theorem}\label{estimate1}
Let \( N \) be a measurable norm on \( H^{1} \). Then, 
\[
\lim _{\varepsilon \to 0} \mathbb{E}\left( \exp \left( \int_{0}^{1} h(s) \, d W_{s} \right) \mid \tilde{N} < \varepsilon \right) = 1,
\]
for all \( h \in L^{2}([0,1]) \).
\end{theorem}

\begin{theorem} \label{result2}  Let $X$ be the solution of \eqref{fir}, and assume that the coefficients satisfy Assumption \((A)\).
For any $\phi \in K_H^\sigma(L^2([0,1]))$ with $\phi_0 = x_0$, when $ \tfrac{1}{2} < H < 1$, 
the Onsager--Machlup functional of $X$ with respect to the norms 
$\Vert \cdot \Vert_\beta$, where $H-\tfrac{1}{2}<\beta<H-\tfrac{1}{4}$, 
 can be expressed as 
	\begin{equation*}
	J(\phi,\dot{\phi})=-\frac{1}{2} \int_0^1 \left(\dot{\phi}_s-s^{\alpha}D_{0^+}^\alpha s^{-\alpha}  \sigma_s^{-1} b_s\left(\phi_s\right)\right)^2+d_H b_x(\phi_s)ds,
\end{equation*}
where
\begin{equation*}
    d_H= \sqrt{\frac{2H\Gamma(H+1/2)\Gamma(3/2-H)}{\Gamma(2-2H)}}.
\end{equation*}
\end{theorem}

\begin{proof}
According to Girsanov theorem, we have
	\begin{align*}
\mathbb{P}(\Vert X-\phi\Vert_\beta \leq \varepsilon)
&=\mathbb{E}\Big(\exp\Big[ -\int_0^1 u_s dW_s - \frac{1}{2} \int_0^1 u_s^2 ds\Big] I_{\Vert \int_0^t \sigma_s B^H_s\Vert_\beta \leq \varepsilon}\Big) \\
  &=	\mathbb{E}\left(\exp(A_1+A_2)I_{\Vert \int_0^t \sigma_s B^H_s\Vert_\beta \leq \varepsilon}\right),
\end{align*}
where
\begin{align*}
	A_1 &= -\int_0^1  \dot{\phi}_s-s^{\alpha}D_{0^+}^\alpha s^{-\alpha}  \sigma_s^{-1} b_s\left(\phi_s+\int_0^s \sigma_u dB^H_u\right)dW_s,\\
 	A_2 &= -\frac{1}{2}\int_0^1 \left(\dot{\phi}_s-s^{\alpha}D_{0^+}^\alpha s^{-\alpha}  \sigma_s^{-1} b_s\left(\phi_s+\int_0^s \sigma_u dB^H_u\right)\right)^2 ds.
\end{align*}
The fractional derivative part therein satisfies
\begin{align*}
	&D_{0^+}^\alpha s^{-\alpha}  \sigma_s^{-1} b_s\left(\phi_s+\int_0^s \sigma_u dB^H_u\right)\\
	=&\frac{1}{\Gamma(1-\alpha)}\left(\frac{b_s(\phi_s+\int_0^s \sigma_u dB^H_u)}{s^{2\alpha}\sigma_s}+\alpha\int_0^s \frac{\frac{b_s(\phi_s+\int_0^s \sigma_udB^H_u)}{s^\alpha \sigma_s}-\frac{b_r(\phi_r+\int_0^r \sigma_u dB^H_u)}{r^\alpha \sigma_r}}{(s-r)^{\alpha+1}}dr \right)\\
	=&\frac{1}{\Gamma(1-\alpha)}\Bigg(\frac{b_s(\phi_s+\int_0^s \sigma_u dB^H_u)}{s^{2\alpha}\sigma_s}+\alpha\int_0^s \frac{\frac{b_s(\phi_s+\int_0^s \sigma_udB^H_u)}{s^\alpha \sigma_s}-\frac{b_r(\phi_s+\int_0^s \sigma_u dB^H_u)}{r^\alpha \sigma_r}}{(s-r)^{\alpha+1}}dr   \\
	&\quad\quad\quad +\alpha\int_0^s \frac{\frac{b_r(\phi_s+\int_0^s \sigma_udB^H_u)}{r^\alpha \sigma_r}-\frac{b_r(\phi_r+\int_0^r \sigma_u dB^H_u)}{r^\alpha \sigma_r}}{(s-r)^{\alpha+1}}dr  \Bigg)\\
	\to & D_{0^+}^\alpha s^{-\alpha}  \sigma_s^{-1}\left(\phi_s' - b_s\left(\phi_s\right)\right),\quad   \Vert  \int_0^\cdot \sigma_s d B^H_s \Vert  \to 0.
\end{align*}


Therefore, the limit of $\mathbb{E}(\exp(A_2)| I_{\Vert \int_0^\cdot \sigma_s dB^H_s\Vert_\beta \leq \varepsilon})$ is 
\begin{equation*}
	\exp\left(-\frac{1}{2}\int_0^1 \left(\dot{\phi}_s-s^{\alpha}D_{0^+}^\alpha s^{-\alpha}  \sigma_s^{-1}b_s\left(\phi_s\right)\right)^2 ds\right).
\end{equation*} 

Next, we decompose $A_1$ into two parts:
\begin{equation*}
	A_1 = B_1 + B_2,
\end{equation*}
where
\begin{align*}
	B_1 &= -\int_0^1 \dot{\phi}_s  dW_s, \\
	B_2 &= \int_0^1 s^{\alpha} D_{0^+}^\alpha s^{-\alpha} \sigma_s^{-1} b_s \left( \phi_s + \int_0^s \sigma_u  dB^H_u \right) dW_s.
\end{align*}
By Theorem~\ref{estimate1}, we obtain
\begin{equation*}
	\lim_{\varepsilon \to 0} \mathbb{E} \left[ \exp(B_1) \,\middle|\, I_{\left\| \int_0^t \sigma_s B^H_s \right\|_\beta \leq \varepsilon} \right] = 1.
\end{equation*}
Applying Taylor's expansion, we further expand $B_2$ as:
\begin{align*}
	B_2 &= \int_0^1 s^{\alpha} D_{0^+}^\alpha s^{-\alpha} \sigma_s^{-1} \left( b_s(\phi_s) + \partial_x b_s(\phi_s) \int_0^s \sigma_u  dB^H_u + R_s \right) ds \\
	    &= C_1 + C_2 + C_3,
\end{align*}
where
\begin{align*}
    C_1 &= \int_0^1 s^{\alpha} D_{0^+}^\alpha s^{-\alpha} \sigma_s^{-1} b_s(\phi_s)  ds, \\
    C_2 &= \int_0^1 s^{\alpha} D_{0^+}^\alpha s^{-\alpha} \sigma_s^{-1} \partial_x b_s(\phi_s) \left( \int_0^s \sigma_u  dB^H_u \right) ds,\\
    C_3 &= \int_0^1 s^{\alpha} D_{0^+}^\alpha s^{-\alpha} \sigma_s^{-1} R_s  ds.
\end{align*}
It is straightforward to show that
\begin{equation*}
	\lim_{\varepsilon \to 0} \mathbb{E} \left[ \exp(C_1) \,\middle|\, I_{\left\| \int_0^\cdot \sigma_s B^H_s \right\|_\beta \leq \varepsilon} \right] = 1.
\end{equation*}
Next, 
\begin{align*}
	C_2
	&=\int_0^1 s^{\alpha}D_{0^+}^\alpha \big(s^{-\alpha} \sigma_s^{-1}\partial_x b_s(\phi_s)\big)\int_0^s \sigma_u\, dB^H_u \, dW_s   \\ %1
	&=  \int_0^1 s^{\alpha}D_{0^+}^\alpha \big(s^{-\alpha} \sigma_s^{-1}\partial_x b_s(\phi_s)\big)\left(\int_0^1 (K_H^*\sigma I_{[0,s]})(u)\, dW_u\right) dW_s \\ %2
	&=\int_0^1 s^{\alpha}D_{0^+}^\alpha \big(s^{-\alpha} \sigma_s^{-1}\partial_x b_s(\phi_s)\big)\left(\int_0^s c_H\Gamma(\alpha)u^{-\alpha}I_{s^-}^\alpha (u^\alpha  \sigma_u)\, dW_u \right)dW_s \\ %3
	&=\int_0^1  \frac{c_H\Gamma(\alpha)}{\Gamma(1-\alpha)} 
	\Bigg(
	 		s^{-\alpha}\sigma_s^{-1}\partial_x b_s(\phi_s)
	 		\int_0^s u^{-\alpha}I_{s^-}^\alpha (u^\alpha  \sigma_u)\, dW_u  \\ %4
	 		&\qquad+\alpha s^\alpha\int_0^s \frac{ 
	 		\big(s^{-\alpha}\sigma_s^{-1}\partial_x b_s(\phi_s)\int_0^s u^{-\alpha}I_{s^{-}}^\alpha (u^\alpha  \sigma_u)\, dW_u\big) 
		}{(s-r)^{\alpha+1}}\,dr  \\ %5
        &\qquad-\alpha s^\alpha\int_0^s \frac{ 
	 		\big(r^{-\alpha}\sigma_r^{-1}\partial_x b_r(\phi_r)\int_0^r u^{-\alpha}I_{r^-}^\alpha (u^\alpha  \sigma_u)\, dW_u\big) 
		}{(s-r)^{\alpha+1}}\,dr \Bigg)\, dW_s \\ %6
		&=\int_0^1 \int_0^s f(s,u)\, dW_u dW_s, %7
\end{align*}
where   
\begin{align*}
	&f(s,u)\\
    =&\frac{c_H\Gamma(\alpha)}{\Gamma(1-\alpha)}	
	\Bigg(
	 		s^{-\alpha}\sigma_s^{-1}\partial_x b_s(\phi_s)u^{-\alpha}I_{s^-}^\alpha (u^\alpha  \sigma_u)  \\
	 		&\quad+\int_0^s \frac{\alpha \sigma_s^{-1}\partial_x b_s(\phi_s)u^{-\alpha}I_{s^-}^\alpha (u^\alpha \sigma_u)}{(s-r)^{\alpha+1}}\,dr
		 -\int_u^s  \frac{\alpha s^\alpha r^{-\alpha} \sigma_r^{-1}\partial_x b_r
         (\phi_r) u^{-\alpha}I_{r^-}^\alpha (u^\alpha \sigma_u) }{(s-r)^{\alpha+1}}\,dr \Bigg) \\
		=:&D_1+D_2+D_3.
\end{align*}

In particular,
\begin{align*}
    D_1 &= \frac{c_H\Gamma(\alpha)}{\Gamma(1-\alpha)} s^{-\alpha}\sigma_s^{-1}\partial_x b(\phi_s)u^{-\alpha}I_{s^-}^\alpha (u^\alpha \sigma_u),   \\
    D_2 &=  \frac{c_H\Gamma(\alpha)}{\Gamma(1-\alpha)} \alpha s^\alpha\int_0^u \frac{ s^{-\alpha}\sigma_s^{-1}\partial_x b(\phi_s)u^{-\alpha}I_{s^-}^\alpha (u^\alpha \sigma_u)}{(s-r)^{\alpha+1}}\,dr, \\
    D_3 &=  \frac{c_H\Gamma(\alpha)}{\Gamma(1-\alpha)} \alpha s^\alpha  \\
    &\quad \int_u^s \frac{ 
		\big(s^{-\alpha}\sigma_s^{-1}\partial_x b_s(\phi_s)u^{-\alpha}I_{s^-}^\alpha (u^\alpha\sigma_u)\big) 
		-\big(r^{-\alpha}\sigma_r^{-1}\partial_x b_r(\phi_r)u^{-\alpha}I_{r^-}^\alpha (u^\alpha\sigma_u)\big)}
		{(s-r)^{\alpha+1}}\,dr . 
\end{align*}

Since 
\begin{align*}
	\lim_{u\to s} I_{s^-}^\alpha (u^\alpha  \sigma_u) 
	&=\lim_{u\to s}\frac{1}{\Gamma(\alpha)}\int_u^s (v-u)^{\alpha-1} v^\alpha  \sigma_v\,dv\\
	&=\lim_{u\to s}\frac{\sigma_s s^\alpha}{\Gamma(\alpha)}\int_u^s (v-u)^{\alpha-1}\,dv\\
	&=\frac{\sigma_s s^\alpha}{\alpha\Gamma(\alpha)}\lim_{u\to s} (s-u)^\alpha\\
	&=0,
\end{align*}
we obtain $D_1(s,s)=0.$

Moreover, as $u\to s$,
\begin{align*}
	\lim_{u\to s}D_2(s,u)
	&= \frac{c_H\Gamma(\alpha)}{\Gamma(1-\alpha)}\lim_{u\to s}\int_0^u \frac{\alpha \sigma_s^{-1}\partial_x b_s(\phi_s)u^{-\alpha}I_{s^-}^\alpha (u^\alpha \sigma_u)}{(s-r)^{\alpha+1}}\,dr\\
	&= \frac{c_H}{\alpha\Gamma(1-\alpha)}\partial_x b(\phi_s),
\end{align*}
hence $D_2(s,s)=\dfrac{c_H}{\alpha\Gamma(1-\alpha)}\partial_x b_s
(\phi_s)$.

Next, decompose $D_3(s,u)$ as
\begin{align*}
	& D_3(s,u)\\ 
    =& \frac{c_H\Gamma(\alpha)}{\Gamma(1-\alpha)} \alpha s^\alpha\int_u^s \frac{ 
	\big(s^{-\alpha}\sigma_s^{-1}\partial_x b_s(\phi_s)u^{-\alpha}I_{s^-}^\alpha (u^\alpha\sigma_u)\big) - \big(r^{-\alpha}\sigma_r^{-1}\partial_x b_r(\phi_r)u^{-\alpha}I_{r^-}^\alpha (u^\alpha\sigma_u)\big)}
		{(s-r)^{\alpha+1}}\,dr \\
		=&: E_1+E_2+E_3,
\end{align*}
where
\begin{align*}
	E_1(s,u)&=\frac{c_H\Gamma(\alpha)}{\Gamma(1-\alpha)} \alpha s^\alpha\int_u^s \frac{ (s^{-\alpha}- r^{-\alpha})\sigma_s^{-1}\partial_x b_s(\phi_s)u^{-\alpha}I_{s^-}^\alpha (u^\alpha \sigma_u)}
		{(s-r)^{\alpha+1}}\,dr,\\
    E_2(s,u)&=\frac{c_H\Gamma(\alpha)}{\Gamma(1-\alpha)} \alpha s^\alpha\int_u^s \frac{ r^{-\alpha}\big(\sigma_s^{-1}\partial_x b_s(\phi_s)-\sigma_r^{-1}\partial_x b_r(\phi_r)\big)u^{-\alpha}I_{s^-}^\alpha (u^\alpha \sigma_u)}
		{(s-r)^{\alpha+1}}\,dr,\\
    E_3(s,u)&=\frac{c_H\Gamma(\alpha)}{\Gamma(1-\alpha)} \alpha s^\alpha\int_u^s \frac{ r^{-\alpha}\sigma_r^{-1}\partial_x b_r(\phi_r)\big(I_{s^-}^\alpha (u^\alpha\sigma_u)- I_{r^-}^\alpha (u^\alpha\sigma_u)\big)}
		{(s-r)^{\alpha+1}}\,dr.
\end{align*}

Observe that
\begin{align*}
    &\lim_{u\to s} \left|E_1(s,u)\right|\\
    =&\frac{c_H\Gamma(\alpha)}{\Gamma(1-\alpha)}\lim_{u\to s}\left| \alpha  \sigma_s^{-1}\partial_x b_s(\phi_s)u^{-\alpha}I_{s^-}^\alpha (u^\alpha\sigma_u)\int_u^s \frac{ s^{-\alpha}- r^{-\alpha}}{(s-r)^{\alpha+1}}\,dr\right|\\
    \leq & \frac{c_H\Gamma(\alpha)}{\Gamma(1-\alpha)}\lim_{u\to s} \left|\alpha  \sigma_s^{-1}\partial_x b_s(\phi_s)u^{-\alpha}\int_u^s (v-u)^{\alpha-1}v^\alpha\sigma_v\,dv\cdot \frac{\alpha}{1-\alpha}u^{-\alpha-1}(s-u)^{1-\alpha}\right| \\
   =&0.
\end{align*}

Moreover, we analyze the limits of \(E_2\) and \(E_3\) as \(u\to s\).

For \(E_2\), we have
\begin{align*}
    |E_2(s,u)|
    &\le \frac{c_H\Gamma(\alpha)}{\Gamma(1-\alpha)}\alpha s^\alpha
    \int_u^s \frac{ r^{-\alpha}\, \big|\sigma_s^{-1}\partial_x b_s(\phi_s)-\sigma_r^{-1}\partial_x b_r(\phi_r)\big|
    \,\big|u^{-\alpha}I_{s^-}^\alpha(u^\alpha\sigma_u)\big|}{(s-r)^{\alpha+1}}\,dr\\
    &\le C\,\sup_{r\in[u,s]}\big|\sigma_s^{-1}\partial_x b_s(\phi_s)-\sigma_r^{-1}\partial_x b_r(\phi_r)\big|
    \int_u^s \frac{(s-u)^\alpha}{(s-r)^{\alpha+1}}\,dr,
\end{align*}
where we used the bound
\[
\big|u^{-\alpha}I_{s^-}^\alpha(u^\alpha\sigma_u)\big|
= \frac{1}{\Gamma(\alpha)}\Big|\int_u^s (v-u)^{\alpha-1} v^\alpha\sigma_v\,dv\Big|
\le C (s-u)^\alpha,
\]
with a constant \(C\) depending on \(\sup_{v\in[0,1]} v^\alpha\sigma_v\). Since
\[
\int_u^s \frac{(s-u)^\alpha}{(s-r)^{\alpha+1}}\,dr = (s-u)^\alpha\int_u^s (s-r)^{-(\alpha+1)}dr
= \frac{(s-u)^\alpha}{\alpha}(s-u)^{-\alpha}=\frac{1}{\alpha},
\]
we obtain the estimate
\[
|E_2(s,u)| \le \frac{C}{\alpha}\,\sup_{r\in[u,s]}\big|\sigma_s^{-1}\partial_x b_s(\phi_s)-\sigma_r^{-1}\partial_x b_r(\phi_r)\big|.
\]
By the continuity of \(\sigma^{-1}\partial_x b_\cdot(\phi_\cdot)\) it follows that
\[
\lim_{u\to s} E_2(s,u)=0.
\]

For \(E_3\) we write the integrand difference of fractional integrals and use continuity of the fractional integral in the upper limit. Indeed,
\begin{align*}
    |E_3(s,u)|
    &\le \frac{c_H\Gamma(\alpha)}{\Gamma(1-\alpha)}\alpha s^\alpha
    \int_u^s \frac{ r^{-\alpha}\,|\sigma_r^{-1}\partial_x b_r(\phi_r)|\,
    \big|I_{s^-}^\alpha(u^\alpha\sigma_u)- I_{r^-}^\alpha(u^\alpha\sigma_u)\big|}{(s-r)^{\alpha+1}}\,dr.
\end{align*}
Using the representation
\[
I_{s^-}^\alpha(u^\alpha\sigma_u)- I_{r^-}^\alpha(u^\alpha\sigma_u)
= \frac{1}{\Gamma(\alpha)}\int_r^s (v-u)^{\alpha-1} v^\alpha\sigma_v\,dv,
\]
we obtain the bound
\[
\big|I_{s^-}^\alpha(u^\alpha\sigma_u)- I_{r^-}^\alpha(u^\alpha\sigma_u)\big|
\le C (s-r)^\alpha,
\]
with \(C\) independent of \(u,r\) (for \(u\) close to \(s\)). Hence
\begin{align*}
    |E_3(s,u)|
    &\le C' \int_u^s \frac{ r^{-\alpha} (s-r)^\alpha}{(s-r)^{\alpha+1}}\,dr
    = C' \int_u^s \frac{r^{-\alpha}}{s-r}\,dr.
\end{align*}
The integral on the right is finite for \(u<s\), and, moreover, by the dominated convergence and continuity of the prefactor \(r^{-\alpha}\sigma_r^{-1}\partial_x b_r(\phi_r)\), we conclude
\[
\lim_{u\to s} E_3(s,u)=0.
\]

Consequently, combining the three pieces \(E_1,E_2,E_3\) we have
\[
\lim_{u\to s} D_3(s,u)=0.
\]

Therefore, gathering previous results,
\[
D_1(s,s)=0,\qquad D_2(s,s)=\frac{c_H}{\alpha\Gamma(1-\alpha)}\partial_x b_s(\phi_s),
\qquad D_3(s,s)=0,
\]
and so
\[
\lim_{u\to s} f(s,u)=\frac{c_H}{\alpha\Gamma(1-\alpha)}\partial_x b(\phi_s).
\]

It follows that  
\begin{align*}
\lim_{\varepsilon \to 0} \mathbb{E} \left[ \exp(C_2) \,\middle|\, \left\| \int_0^t \sigma_s\, dB^H_s \right\|_\beta \leq \varepsilon \right] 
&= \exp\left( -\frac{1}{2} \int_0^1 f(s,s)\,ds \right) \\
&= \exp\left( -\frac{d_H}{2} \int_0^1 \partial_x b_s(\phi_s)\,ds \right),
\end{align*}
where
\begin{equation*}
    d_H= \sqrt{\frac{2H\Gamma(H+1/2)\Gamma(3/2-H)}{\Gamma(2-2H)}}.
\end{equation*}










Finally, we analyze the limiting behavior of the term $C_3$. Since 
\begin{align*}
    R_s&=b_s\left(\phi_s+\int_0^s \sigma_udB^H_u\right)-b_s(\phi_s)-b_s(\phi_s)\int_0^s \sigma_udB^H_u\\
    &=\int_0^1 \int_0^\lambda \partial_{xx}b_s(\mu\int_0^s \sigma_udB^H_u+\phi_s) \left(\int_0^s \sigma_udB^H_u\right)^2 d\mu d\lambda.
\end{align*}
and
\begin{align*}
&|R_s - R_r| \\
\leq &
  \bigg(\int_0^s \sigma_u\, dB^H_u\bigg)^2 
  \int_0^1 \int_0^\lambda 
  \Big|
     \partial_{xx}b_s\!\left(\mu \int_0^s \sigma_u\, dB^H_u + \phi_s\right)\\
     &- \partial_{xx}b_r\!\left(\mu \int_0^r \sigma_u\, dB^H_u + \phi_r\right)
  \Big| \, d\mu\, d\lambda \\
& + 
  \Bigg|
     \bigg(\int_0^s \sigma_u\, dB^H_u\bigg)^2 
     - \bigg(\int_0^r \sigma_u\, dB^H_u\bigg)^2
  \Bigg| 
  \int_0^1 \int_0^\lambda 
     \Big|\partial_{xx}b_s\!\left(\mu \int_0^s \sigma_u\, dB^H_u + \phi_s\right)\Big| 
  \, d\mu\, d\lambda \\[1ex]
\leq  &
  \int_0^1 \int_0^\lambda 
     L_3 \bigg(\int_0^s \sigma_u\, dB^H_u\bigg)^2 
     \Big(
        \mu\Big|\int_0^s \sigma_u\, dB^H_u - \int_0^r \sigma_u\, dB^H_u\Big|
        + |\phi_s - \phi_r|
     \Big)\, d\mu\, d\lambda \\
& +
  \int_0^1 \int_0^\lambda 
     L_3 \bigg(\int_0^s \sigma_u\, dB^H_u\bigg)^2 
     \Big(
        2\|b\|_\infty |s-r| + 2|s-r|^H
     \Big)\, d\mu\, d\lambda \\
& +
  \int_0^1 \int_0^\lambda 
     \mu \|b_{xx}\|_\infty 
     \Bigg|
        \bigg(\int_0^s \sigma_u\, dB^H_u\bigg)^2 
        - \bigg(\int_0^r \sigma_u\, dB^H_u\bigg)^2
     \Bigg| \, d\mu\, d\lambda \\[1ex]
\leq &
  L_3 \Big(
        |\phi_s - \phi_r| 
        + \Big|\int_0^s \sigma_u\, dB^H_u - \int_0^r \sigma_u\, dB^H_u\Big|
        + 2\|b\|_\infty |s-r|
        + 2|s-r|^H
     \Big)
     \bigg(\int_0^s \sigma_u\, dB^H_u\bigg)^2 \\
& + 
  \|b_{xx}\|_\infty 
  \Bigg|
     \bigg(\int_0^s \sigma_u\, dB^H_u\bigg)^2 
     - \bigg(\int_0^r \sigma_u\, dB^H_u\bigg)^2
  \Bigg|,
\end{align*}
we obtain
\begin{align*}
     & |s^{\alpha}D_{0^+}^\alpha s^{-\alpha} \sigma_s^{-1}R_s|\\
     =&\frac{1}{\Gamma(1-\alpha)}  s^{\alpha}\left|\left( \frac{ s^{-\alpha} \sigma_s^{-1}R_s}{s^\alpha} + \alpha \int_0^s \frac{ s^{-\alpha} \sigma_s^{-1}R_s -  r^{-\alpha} \sigma_r^{-1}R_r}{(s - r)^{\alpha+1}}  dr \right)\right|\\
     \leq & C\left| \frac{\varepsilon^2}{s^\alpha}+s^{-\alpha}\int_0^1 (1-x^{-\alpha})(1-x)^{\alpha+1}dx  \right. \\
     &\quad + \left.  \int_0^s \frac{r^{-\alpha} R_s(\sigma_s^{-1}-\sigma_r^{-1})}{(s-r)^{\alpha+1}}dr+\int_0^s \frac{r^{-\alpha}\sigma_r^{-1}(R_s-R_r)}{(s-r)^{\alpha+1}}ds   \right|\\
     \leq & C \varepsilon^2.
\end{align*}



By proceeding in the same way as in the case $H<\tfrac{1}{2}$, we obtain 
\[
\lim_{\varepsilon \to 0} \mathbb{E} \left[ \exp(C_3) \,\middle|\, \left\| \int_0^t \sigma_s\, dB^H_s \right\|_\beta \leq \varepsilon \right] = 1.
\]




Combining the above results, the functional takes the form  
\[
J(\phi,\dot{\phi}) = -\frac{1}{2} \int_0^1 \left(\dot{\phi}_s -s^{\alpha} D_{0^+}^\alpha s^{-\alpha} \sigma_s^{-1}  b_s(\phi_s)  \right)^2 +d_H\partial_x b_s(\phi_s)\, ds .
\]




\end{proof}

\end{document}
